\SourceLeaf{309}%
giving the method by which this expansion can be carried as far as desired.

I take the liberty of making a few changes to the notation used there, partly to make printing easier and partly to give the notation in the various parts of the present memoir a little more symmetry. Here, to begin with, let
\begin{list}{}{\setlength{\leftmargin}{1.4em}\setlength{\itemindent}{-1.4em}\setlength{\itemsep}{0pt}\setlength{\parsep}{0pt}}
\item \(r\) denote the distance between the two points, taking the radius of the sphere as the unit.
\item \(B+\frac12b\) and \(B-\frac12b\) denote the latitudes at the first and second endpoints of \(r\).
\item \(T+\frac12t\) and \(T-\frac12t\pm180^0\) denote the azimuth of the second endpoint from the first and of the first endpoint from the second, respectively.
\item \(l\) denote the difference in longitude.
\end{list}
\noindent We assume that the azimuth is counted from south toward west and that \(l\) is positive when the second point lies farther west than the first.

We also set
\[
\begin{aligned}
\beta&=r\cos T\\
\tau&=r\sin T\mathord{.}\operatorname{tang}B\\
\lambda&=\frac{r\sin T}{\cos B}
\end{aligned}
\]
These quantities express what was denoted by \(b^0,a^0,\lambda^0\) in Art. 17, namely the values of \(b,t,l\) accurate up to but excluding third order. They satisfy the equation
\[
rr+\tau\tau=\beta\beta+\lambda\lambda
\]
Throughout, the orders are understood by treating \(r\) as a quantity of first order.

\indent For brevity we also write
\[
\begin{aligned}
\frac{2\operatorname{tang}\frac12r}{r}&=m\\
\frac{2\sin\frac12r}{r}&=n
\end{aligned}
\]
The easiest way to obtain the intended expansion is to convert the equation
\[
x=\sin y
\]
